\section{Zadanie 14}

Udowodnij, że dla każdego dodatniego $n$ istnieje turniej (graf skierowany
pełny) na $n$ wierzchołkach, który ma co najmniej $\frac{n!}{2^{n-1}}$ ścieżek Hamiltona.

\subsection*{Rozwiązanie}

Rozważmy nieskierowany graf pełny na $n$ wierzchołkach. Będziemy kierować wszystkie jego krawędzie z równym prawpodpodobieństwem dla dwóch stron. Policzmy wartość oczekiwaną liczby ścieżek Hamiltona dla naszego grafu. Z liniowośći wartości oczekiwanej możemy rozbić ją na sumę wartości oczekiwanych indykatorów czy dana permutacja wierzchołków jest ścieżką Hamiltona. Permutacji jest $n!$, dla każdej z nich wartość oczekiwana jest równa prawdopodobieństwu i wynosi $\frac{1}{2^{n-1}}$ (mamy $n-1$ krawędzi, tylko jedno ze skierowań daje ścieżkę Hamiltona). Po zsumowaniu dostajemy wartość oczekiwaną ścieżek Hamiltona w grafie $\frac{n!}{2^{n-1}}$, z czego na mocy uwagi z zadania 13 wynika, że prawdopodbieństwo tego, że graf ma conajmniej tyle ścieżek Hamiltona jest większe od zera, czyli musi istnieć takie skierowanie które spełnia warunki zadania. 